% VI34-DETAILED-PROBLEM-11
\begin{problem}[VI.34.P11: Degree-zero scaling invariance]\label{prob:vi34-11}
Let \(F/G\) be a homogeneous rational expression. Explain algebraically why equal numerator and denominator degree is the scaling-invariant condition.
\end{problem}

\ifdefined\FullProblemDossiers
\begin{solution}
Suppose \(F\) and \(G\) are homogeneous of degrees \(a\) and \(b\). Under simultaneous scaling,
\[
F(\lambda x)=\lambda^aF(x),
\qquad
G(\lambda x)=\lambda^bG(x).
\]
Therefore
\[
\frac{F(\lambda x)}{G(\lambda x)}
=
\lambda^{a-b}\frac{F(x)}{G(x)}.
\]
If \(a=b\), the factor is \(1\), so the ratio is invariant.

The intrinsic graded statement avoids exceptional behavior caused by evaluating only on the finite set \(k^\times\):
the rational expression has graded degree
\[
a-b,
\]
and it belongs to the degree-zero part exactly when
\[
\boxed{a=b}.
\]
Thus degree-zero localization is the algebraic invariant notion, independently of the cardinality of the ground
field.
\end{solution}
\fi
